\chapter*{Conclusion générale}
\addcontentsline{toc}{chapter}{Conclusion générale}
\markboth{Conclusion générale}{}

Le recrutement confronte les organisations à un volume important de CV, à des offres exprimées de manière hétérogène et à la nécessité de justifier un premier tri. Dans ce contexte, rapprocher automatiquement un CV et une offre ne doit pas réduire une candidature à un classement opaque. Le besoin auquel répond JobTest est celui d'une aide structurée~: mettre en évidence les éléments de compatibilité, les écarts et les raisons d'une recommandation, tout en laissant au recruteur la décision finale.

Le projet avait pour objectif de concevoir et réaliser une plateforme Web d'aide au recrutement capable de gérer les offres, les candidats et leurs CV, de produire une évaluation hybride, d'en restituer une explication et d'en préserver la traçabilité. La solution livrée propose un frontend Next.js, une API FastAPI, une persistance SQLAlchemy utilisant SQLite en développement et PostgreSQL dans l'environnement conteneurisé, ainsi qu'un export PDF. Le recruteur peut créer ses offres, gérer ses candidats, lancer une évaluation, consulter son historique et obtenir un rapport. Les contrôles d'appartenance des données maintiennent la séparation entre les espaces des recruteurs.

La démarche a suivi les étapes de l'étude du contexte, de l'analyse des besoins, de la conception UML et architecturale, du développement, de la conception du moteur de rapprochement, puis de la validation et du déploiement. Elle s'inscrit dans le cycle en cascade présenté au début du rapport, tout en intégrant les retours techniques nécessaires au fil des vérifications. L'industrialisation a complété le développement par Docker Compose, un proxy inverse comme point d'entrée et un déploiement de démonstration sur une VM Microsoft Azure, servi en HTTPS.

La contribution centrale est un moteur hybride qui combine les signaux TF--IDF, sémantique et compétences, contrôle les preuves présentes dans le CV, calibre son score puis applique des règles de décision explicites. Le score est reproductible, versionné et s'accompagne de sous-scores, de forces, de lacunes et de recommandations structurées. Un CV qui recopie les exigences de l'offre ou qui ne fait que les lister n'obtient pas le verdict \texttt{HIRE}. Une seconde couche locale, exécutée avec Ollama, produit une analyse qualitative séparée~: elle enrichit la lecture du recruteur, mais ne peut ni modifier le score officiel ni déterminer seule la recommandation algorithmique \texttt{HIRE}, \texttt{CONSIDER} ou \texttt{REJECT}. Cette séparation est un résultat important du projet, car elle maintient une distinction claire entre calcul déterministe et contenu génératif.

Les validations réunies montrent que le parcours complet fonctionne~: authentification, création d'offre, évaluation d'un texte ou d'un document, persistance, analyse approfondie et génération PDF. Je l'ai exercé sur des données synthétiques, sur des CV fictifs au format PDF et sur mes propres CV face à des offres réelles. Le déploiement a également confirmé les cinq services conteneurisés, le proxy inverse, l'accès HTTPS et l'isolation des ports métier du réseau public. Sur le test synthétique isolé, les verdicts affichés par la version finale atteignent un macro-F1 de 0,740 sur trois classes, sans dépasser TF--IDF seul calibré. Ces mesures restent internes au corpus synthétique~; les essais menés avec mes propres CV, faute de jugements de référence indépendants, restent qualitatifs. La valeur démontrée est celle d'un prototype fonctionnel et traçable, non celle d'un système déjà validé en situation opérationnelle.

Cette prudence est essentielle. Le jeu d'évaluation quantitative est synthétique, les ancrages de calibration sont appris sur ses étiquettes, les poids restent fixés empiriquement, le dictionnaire de compétences est nécessairement incomplet et aucune annotation par des recruteurs n'a été réalisée. Un CV qui décrit faussement une expérience n'est pas détecté par l'analyse du texte. Les biais possibles dans les données, les libellés ou les règles doivent être étudiés avant tout usage réel. La couche 2 présente une latence notable sur CPU et son texte, bien que contrôlé par des statuts, reste consultatif. Du point de vue de l'infrastructure, la démonstration est hébergée sur une VM unique, sans haute disponibilité, sans sauvegarde automatisée ni supervision avancée. Une politique explicite de protection, de conservation, d'anonymisation et de suppression des CV demeure également indispensable.

Les évolutions prioritaires à court terme sont la sécurisation des sessions par cookies adaptés, la restriction de SSH, les sauvegardes PostgreSQL, le monitoring et la correction des limites de téléversement. L'amélioration de l'interface et l'optimisation de la couche 2 peuvent progresser en parallèle. À moyen terme, il sera nécessaire de confronter le système à un corpus réel anonymisé et autorisé, de travailler avec des recruteurs, d'enrichir le dictionnaire de compétences, de réapprendre la calibration du score sur ces données annotées, ce qui constitue la priorité, d'évaluer les biais et de formaliser l'évaluation des explications génératives, y compris dans plusieurs langues.

À plus long terme, une charge réelle justifierait des traitements asynchrones, une file de tâches, une meilleure observabilité et, si nécessaire, la séparation de certains services sur une infrastructure évolutive. Des intégrations avec des systèmes de gestion de recrutement pourraient être envisagées, accompagnées de mécanismes de gouvernance et d'audit des décisions assistées. Ces perspectives ne sont pas des fonctions déjà réalisées~; elles dessinent les conditions d'une évolution responsable.

JobTest constitue ainsi un prototype fonctionnel et une base technique cohérente pour l'aide à l'analyse des candidatures. Il organise des signaux et des explications utiles au recruteur, sans prétendre remplacer son jugement, l'appréciation du contexte ou sa responsabilité dans la décision finale.
